\section{USAMO 1995}
        \begin{enumerate}
        	\item (\textit{USAMO 1995}) Let $\, p \,$ be an odd \href{https://artofproblemsolving.com/wiki/index.php?title=Prime}{prime}.  The sequence $(a_n)_{n \geq 0}$ is defined as follows: $\, a_0 = 0,$ $a_1 = 1, \, \ldots, \, a_{p-2} = p-2 \,$ and, for all $\, n \geq p-1, \,$ $\, a_n \,$ is the least positive integer that does not form an arithmetic sequence of length $\, p \,$ with any of the preceding terms. Prove that, for all $\, n, \,$ $\, a_n \,$ is the number obtained by writing $\, n \,$ in base $\, p-1 \,$ and reading the result in base $\, p$.


 



\textbf{Solution}

I have to make the assumption that $a_{n+1}>a_{n}$ (without this assumption, I can have the sequence

 ${1,\cdots,p-2,1,1,\cdots,1,2,1,\cdots}$)

 All of the following work are in base $p$ otherwise stated

 


 Lemma 1: for an arithmetic sequence of length $\, p \,$ to exist, there must be a number in the sequence with $(p-1)$ as a digit.

 A arithmetic sequence of length $\, p \,$ can be represented as

 $a,a+d,a+2d,\cdots,a+(p-1)d$

 Since no number repeats, $d\neq0$. Thus, d must have a rightmost non-zero digits, and every term in the sequence $0,d,2d,\cdots,(p-1)d$ have the same number of tailing zeros, let's say there are $x$ tailing zeros.

 and let $\dfrac{{0,d,2d,\cdots,(p-1)d}}{p^x}=S$

 $S\equiv{1,2,...,(p-1)}(mod)p$ since p is an odd prime and the operation divide by $p^x$ has remove all factors of p in S

 Thus, there must be a number with $(p-1)$ as a digit if a length-p sequence exist.

 


 Now, I'm going to prove the statement by strong induction.

 I'm going to assume that for all $\, k \,$  less than $\, n \,$ 

 $\, a_k \,$ is the number obtained by writing $\, k \,$ in base $\, p-1 \,$ and reading the result in base $\, p$.

 (which is true for $n=p-2$ already)

 Or in another word, the terms precede $a_n$ contains all number less than $n$ written in base $\, p-1 \,$ and reading the result in base $\, p$.

 


 Lemma 2: Any term containing the digit $p-1$ will form an arithmetic sequence of length-p with preceding terms

 let $p-1$ be the $x^{th}$ digit from the right (There may be more than 1 $(p-1)$)

 $a=$ the number with all $(p-1)$ replaced by 0 (e.g: $p=7$, the number$=1361264$, then $a=1301204$)

 $d=\sum p^{x-1}$ (for all $x$'s)

 $a,a+d,a+2d,\cdots,a+(p-2)d$ all precede $a_n$, thus, $a+(p-1)d$ can't be in the sequence $(a_n)_{n \geq 0}$

 Therefore, any term containing the digit $p-1$ can't be a term in the sequence ${(a_n)}_{n\ge0}$

 


 Since the digit $p-1$ won't appear (lemma 2) and as long as it doesn't appear, the arithmetic sequence of length $\, p \,$ won't be formed (lemma 1), and $a_n$ must be as small as possible, 

 Therefore, for all $\, n, \,$ $\, a_n \,$ is the number obtained by writing $\, n \,$ in base $\, p-1 \,$ and reading the result in base $\, p$.

 


	\item (\textit{USAMO 1995}) A calculator is broken so that the only keys that still work are the $\, \sin, \; \cos,$ $\tan, \; \sin^{-1}, \; \cos^{-1}, \,$ and $\, \tan^{-1} \,$ buttons.  The display initially shows 0. Given any positive rational number $\, q, \,$ show that pressing some finite sequence of buttons will yield $\, q$.   Assume that the calculator does real number calculations with infinite precision.  All functions are in terms of radians.


 



\textbf{Solution}

We will prove the following, stronger statement :  If $m$ and $n$ are relatively prime nonnegative integers such that $n>0$, then the some finite sequence of buttons will yield $\sqrt{m/n}$.

 To prove this statement, we induct strongly on $m+n$.  For our base case, $m+n=1$, we have $n=1$ and $m=0$, and $\sqrt{m/n} = 0$, which is initially shown on the screen.  For the inductive step, we consider separately the cases $m=0$, $0<m\le n$, and $n<m$.

 If $m=0$, then $n=1$, and we have the base case.

 If $0< m \le n$, then by inductive hypothesis, $\sqrt{(n-m)/m}$ can be obtained in finitely many steps; then so can
 \[\cos \tan^{-1} \sqrt{(n-m)/m} = \sqrt{m/n} .\]

 If $n<m$, then by the previous case, $\sqrt{n/m}$ can be obtained in finitely many steps.  Since $\cos \tan^{-1} \sqrt{n/m} = \sin \tan^{-1} \sqrt{m/n}$, it follows that
 \[\tan \sin^{-1} \cos \tan^{-1} \sqrt{n/m} = \sqrt{m/n}\]can be obtained in finitely many steps.  Thus the induction is complete.  $\blacksquare$

 


 \textbf{Alternate Solution (using infinite descent):} This solution is similar to the above solution. We can invert a number with
 \[\tan \sin^{-1} \cos \tan^{-1} \sqrt{n/m} = \sqrt{m/n}\] and we can get 
 \[\tan \sin^{-1} \cos \tan^{-1}\cos\tan^{-1}\sqrt{x} = \sqrt{x + 1}.\]Then we do $\sqrt{x}$ to $\sqrt{x + 1}$ in reverse, trying to get 1. Subtracting one inside the radical for a fraction $\sqrt{m/n}$ will eventually result in a fraction less than or equal to 1. Taking the reciprocal of this will result in a fraction greater than or equal to 1, and we can keep subtracting one inside the radical.

 We do this repeatedly, and this should eventually result in 1 in a finite number of steps. This is because each time we do this, m+n is smaller. This is true even if $\sqrt{m/n}$ can be simplified, as simplifying will only result in a smaller value. The smallest possible value of the fraction is 1/1, so this will be achieved after finite steps.

 Finally, we know that 1 is achievable as cos(0) = 1.

 
<i>Alternate solutions are always welcome. If you have a different, elegant solution to this problem, please \href{https://artofproblemsolving.com/wiki/index.php?title=1995\_USAMO\_Problems/Problem\_2&action=edit}{add it} to this page.</i>

 


	\item (\textit{USAMO 1995}) Given a nonisosceles, nonright triangle $\, ABC, \,$ let $\, O \,$ denote the center of its circumscribed circle, and let $\, A_1, \, B_1, \,$ and $\, C_1 \,$ be the midpoints of sides $\, BC, \, CA, \,$ and $\, AB, \,$ respectively.  Point $\, A_2 \,$ is located on the ray $\, OA_1 \,$ so that $\, \triangle OAA_1 \,$ is similar to $\, \triangle OA_2A$.  Points $\, B_2 \,$ and $\, C_2 \,$ on rays $\, OB_1 \,$ and $\, OC_1, \,$ respectively, are defined similarly.  Prove that lines $\, AA_2, \, BB_2, \,$ and $\, CC_2 \,$ are concurrent, i.e. these three lines intersect at a point.


 



\textbf{Solution 1}

\textbf{LEMMA 1: } In $\triangle ABC$ with circumcenter $O$, $\angle OAC = 90 - \angle B$. 

 \textbf{PROOF of Lemma 1:} The arc $AC$ equals $2\angle B$ which equals $\angle AOC$.  Since $\triangle AOC$ is isosceles we have that $\angle OAC = \angle OCA = 90 - \angle B$. QED

 Define $H \in BC$ s.t. $AH \perp BC$.  Since $OA_1 \perp BC$, $AH \parallel OA_1$.  Let $\angle AA_2O = \angle A_1AO = x$ and $\angle AA_1O = \angle A_2AO = y$.  Since we have $AH \parallel OA_1$, we have that $\angle HAA_2 = x$.  Also, we have that $\angle A_2AA_1 = y-x$.  Furthermore, $\angle BAH = 90 - \angle B = \angle OAC$, by lemma 1.  Therefore, $\angle A_1AC = 90 - \angle B + x = \angle BAA_2$.  Since $A_1$ is the midpoint of $BC$, $AA_1$ is the median.  However $\angle A_1AC =  \angle BAA_2$ tells us that $AA_2$ is just $AA_1$ reflected across the internal angle bisector of $A$.  By definition, $AA_2$ is the $A$-symmedian.  Likewise, $BB_2$ is the $B$-symmedian and $CC_2$ is the $C$-symmedian.  Since the symmedians concur at the symmedian point, we are done.

 QED

 



\textbf{Solution 2}

Let $AH$ be the altitude of $\triangle ABC \implies$
 \[\angle BAH = 90^\circ - \angle ABC, \angle OAC  = \angle OCA =\]
 \[=  \frac{180^\circ - \angle AOC}{2} = \frac{180^\circ - 2\angle ABC}{2} = \angle BAH.\]Hence $AH$ and $AO$ are isogonals with respect to the angle $\angle BAC.$
 \[\triangle OAA_1 \sim \triangle OA_2A, AH || A_1O \implies \angle AA_1O = \angle A_2AO = \angle A_1AH.\]$AA_2$ and $AA_1$ are isogonals with respect to the angle $\angle BAC.$

 Similarly $BB_2$ and $BB_1$ are isogonals with respect to $\angle ABC.$

 Similarly $CC_2$ and $CC_1$ are isogonals with respect to $\angle ACB.$

 Let $G = AA_1 \cap BB_1$ be the centroid of $\triangle ABC, L = AA_2 \cap BB_2.$

 $L$ is the isogonal conjugate of a point $G$ with respect to a triangle $\triangle ABC.$

 
 \[G \in CC_1 \implies L \in CC_2.\]

 <i>\textbf{Corollary}</i>

 If median and symmedian start from any vertex of the triangle, then the angle formed by the symmedian and the angle side has the same measure as the angle between the median and the other side of the angle.

 $AA_1, BB_1, CC_1$ are medians, therefore $AA_2, BB_2, CC_2$ are symmedians, so the three symmedians meet at a point which is triangle center called the Lemoine point.

 \textbf{vladimir.shelomovskii@gmail.com, vvsss}

 


	\item (\textit{USAMO 1995}) Suppose $\, q_0, \, q_1, \,  q_2, \ldots \; \,$ is an infinite sequence of integers satisfying the following two conditions:

(i)  $\, m-n \,$ divides $\, q_m - q_n \,$ for $\, m > n \geq 0,$ 

(ii) there is a polynomial $\, P \,$ such that $\, |q_n| < P(n) \,$ for all $\, n$. 

Prove that there is a polynomial $\, Q \,$ such that $\, q_n = Q(n) \,$ for all $\, n$.


 



\textbf{Solution}

Step 1: Suppose $P$ has degree $d$. Let $Q$ be the polynomial of degree at most $d$ with $Q(x)=q_x$ for $0\leq x\leq d$. Since the $q_x$ are all integers, $Q$ has rational coefficients, and there exists $k$ so that $kQ$ has integer coefficients. Then $m-n|kQ(m)-kQ(n)$ for all $m,n\in \mathbb N_0$.

 Step 2: We show that $Q$ is the desired polynomial.

 Let $x>n$ be given. Now
 \[kq_x\equiv kq_m\pmod{x-m}\text{ for all integers }m\in[0,d]\]Since $kQ(x)$ satisfies these relations as well, and $kq_m=kQ(m)$,
 \[kq_x\equiv kQ(x)\pmod{x-m}\text{ for all integers }m\in[0,d]\]and hence
 \[kq_x\equiv kQ(x)\pmod{\text{lcm}(x,x-1,\ldots, x-d)}. \;(1)\]Now
 \begin{align*} \text{lcm}(x,x-1,\ldots, x-i-1)&=\text{lcm}[\text{lcm}(x,x-1,\ldots, x-i),x-i-1]\\ &=\frac{\text{lcm}(x,x-1,\ldots, x-i)(x-i-1)}{\text{gcd}[\text{lcm}(x,x-1,\ldots, x-i),(x-i-1)]}\\ &\geq \frac{\text{lcm}(x,x-1,\ldots, x-i)(x-i-1)}{\text{gcd}[x(x-1)\cdots(x-i),(x-i-1)]}\\ &\geq \frac{\text{lcm}(x,x-1,\ldots, x-i)(x-i-1)}{(i+1)!} \end{align*}
so by induction $\text{lcm}(x,x-1,\ldots, x-d)\geq \frac{x(x-1)\cdots (x-d)}{d!(d-1)!\cdots 1!}$. Since $P(x), Q(x)$ have degree $d$, for large enough $x$ (say $x>L$) we have $\left|Q(x)\pm\frac{x(x-1)\cdots (x-d)}{kd!(d-1)!\cdots 1!}\right|>P(x)$. By (1) $kq_x$ must differ by a multiple of $\text{lcm}(x,x-1,\ldots, x-d)$ from $kQ(x)$; hence $q_x$ must differ by a multiple of $\frac{x(x-1)\cdots (x-d)}{kd!(d-1)!\cdots 1!}$ from $Q(x)$, and for $x>L$ we must have $q_x=Q(x)$.

 Now for any $y$ we have $kQ(y)\equiv kQ(x)\equiv kq_x \equiv kq_y\pmod{x-y}$ for any $x>L$. Since $x-y$ can be arbitrarily large, we must have $Q(y)=q_y$, as needed.

 


	\item (\textit{USAMO 1995}) Suppose that in a certain society, each pair of persons can be classified as either \textit{amicable} or \textit{hostile}. We shall say that each member of an amicable pair is a \textit{friend} of the other, and each member of a hostile pair is a \textit{foe} of the other.  Suppose that the society has $\, n \,$ persons and $\, q \,$ amicable pairs, and that for every set of three persons, at least one pair is hostile.  Prove that there is at least one member of the society whose foes include $\, q(1 - 4q/n^2) \,$ or fewer amicable pairs.  


 



\textbf{Solution}

Consider the graph with two people joined if they are friends.

 For each person $X$, let $A(X)$ be the set of its friends and $B(X)$ the set of its foes. Note that any edge goes either: from $X$ to $A(X)$ (type 1), from $A(X)$ to $B(X)$ (type 2) or from a point of $B(X)$ to another (type 3), but there's no edge joining two points of $A(X)$ (since they would form a triangle with $X$). Let the number of type 1, type 2, type 3 edges of $X$ be $x_1, x_2, x_3$ respectively, so that $x_1$ is the degree of $X$ and we want to show that for some $X$, we have $x_3 \leq q(1-4q/n^2)$.

 Since each edge is of one of those types, we have $x_1+x_2+x_3=q$. Thus
 \[x_3 \leq q(1-4q/n^2) \Longleftrightarrow \\ x_3 \leq q-4q^2/n^2 \Longleftrightarrow \\ q-x_1-x_2 \leq q-4q^2/n^2 \Longleftrightarrow \\ x_1+x_2 \geq 4q^2/n^2.\]That is, what we want is equivalent to proving that for some vertex $X$, the set of edges touching either $X$ or a vertex joined to $X$ is at least $4q^2/n^2$. Obviously now we'll sum $x_1+x_2$ over all vertices $X$. In the resulting sum, an edge joining $X, Y$ is counted once for each edge that $X, Y$ have, that is, it is counted $D(X)+D(Y)$ times, where $D(X)=x_1$ is the degree of $X$. Thus each vertex $X$ contributes to the overall sum with $D(X)$ for each edge it has, and since it has $D(X)$ edges, it contributes with $D(X)^2$. Thus the considered sum is equal to
 \[\sum_{X \in G} D(X)^2 \geq \frac{(\sum_{X \in G} D(X))^2}{n}=\frac{(2q)^2}{n}=4q^2/n.\](That's Cauchy.) Since we are summing over $n$ vertices, one of the summands is at least $4q^2/n^2$ by pigeonhole, which is what we wanted to prove.

 <i>Alternate solutions are always welcome. If you have a different, elegant solution to this problem, please \href{https://artofproblemsolving.com/wiki/index.php?title=1995\_USAMO\_Problems/Problem\_5&action=edit}{add it} to this page.</i>

 


\end{enumerate}